\chapter{Einleitung und Fragestellung}
\label{chap:einleitung}

Die vorliegende Dokumentation beschreibt die Durchführung und die Ergebnisse unseres Datenanalyseprojekts im Rahmen des Moduls Data-Analytics.
Gemäß der im Kurs behandelten Phasen liegt der Fokus auf der Definition einer Fragestellung,
der methodischen Datenbeschaffung sowie der Anwendung moderner statistischer Verfahren zur Ableitung von Handlungsempfehlungen.

\section*{Armut im viktorianischen London}
\label{sec:kontext}
Während des 19.\ Jahrhunderts versuchte der Positivist Charles Booth, Ausmaß und Ursachen der Armut in East-London zu erfassen.
Dafür nutzte er bestehende statistische Daten,
die er dann analysierte, kategorisierte, und schließlich 1897 in \textit{"LIFE AND LABOUR OF THE PEOPLE IN LONDON"}\ (im Folgenden auch als \textit{"LIFE AND LABOUR"} abgekürzt)
\cite{life_and_labour} veröffentlichte,
welches eines der bekanntesten Exemplare einer frühen (vor-) soziologischen Studie ist.

Dafür führt er die in \autoref{tab:booth_classes} beschriebene Klassifikation der Arbeiterklasse ein.

\begin{table}[htbp]
\centering
\caption{Charles Booths Klassifikation sozialer Schichten}
\label{tab:booth_classes}

\begin{tabular}{|c|p{4.8cm}|c|p{3.8cm}|}
\hline
\textbf{Class} & \textbf{Description} & \textbf{Category} & \textbf{Alias} \\
\hline

A & Lowest class &
very poor \& semi-criminal &
Lowest Class \\
\hline

B & Casual earnings &
very poor \& semi-criminal &
Casual Earnings \\
\hline

C & Intermittent earnings &
poor &
Irregular Earnings \\
\hline

D & Small regular earnings &
poor &
Regular Minimum \\
\hline

E & Regular standard earnings &
comfortable &
Ordinary Earnings \\
\hline

F & Higher class labour &
comfortable &
Highly Paid Work \\
\hline

G & Lower middle class &
well to do &
Lower Middle \\
\hline

H & Upper middle class &
well to do &
Upper Middle \\
\hline

\end{tabular}
\end{table}

\section*{Motivation und Relevanz}
In der Entstehungsgeschichte der Soziologie sind Armut und Ungleichheit neben Religion, und Angewohnheiten und Verhalten zentrale Themen.
Es ging dabei häufig darum, den Ausmaß von Armut im Zuge der Industriellen Revolution und seiner Folgen zu Messen.
Auch in modernen Gesellschaften sind Armut und Ungleichheit weit verbreitet, wenn auch nicht im selben Ausmaß.
Nichtsdestotrotz bietet der Blick auf die Zeit, in der unsere modernen Gesellschaften entstanden eine wertvolle Perspektive.

Dafür werden in dieser Arbeit mögliche Faktoren für Armut und die Entstehung sozialer Ungleichheit untersucht.

\subsection*{Methodische Motivation: Moderne Analyse historischer Daten}
Da Charles Booth in seiner Zeit keinen Zugang zu modernen statistischen Methoden und Tests hatte, ist auch interessant,
inwiefern wir mit diesen neue Einsichten erlangen können.
Insbesondere ist auch eine grafische Aufbereitung der Daten, die über die Bar-Charts in \textit{"LIFE AND LABOUR"} hinausgeht möglich.

\section*{Verfeinerung der Fragestellung}
\label{sec:verfeinerung}

Unsere ursprüngliche Fragestellung lautet: \textit{„Welche Faktoren bestimmten die Armut im London des 19. Jahrhunderts?“}.

\subsection*{Konkretisierung und Hypothesenbildung}
Um die Fragestellung operabel zu machen, haben wir sie in drei konkrete Untersuchungsebene unterteilt:

\begin{figure}[h]
    \centering
    \begin{tikzpicture}[
        scale=0.9,
        transform shape,
        scale=0.8,
        transform shape,
        node distance=2.8cm and 2.8cm,
        every node/.style={font=\small},
        main/.style={
            rectangle,
            rounded corners,
            draw=black,
            thick,
            align=center,
            minimum width=4.8cm,
            minimum height=1.2cm,
            fill=gray!10
        },
        sub/.style={
            rectangle,
            rounded corners,
            draw=black,
            align=center,
            text width=4.8cm,
            minimum height=3.6cm,
            fill=blue!5
        },
        arrow/.style={
            ->,
            thick
        }
    ]

    % main node
    \node[main] (poverty)
    {\textbf{Armut im viktorianischen London}};

    % sub node
    \node[sub, below left=of poverty] (demo)
    {
        \textbf{Demografisches Risiko}\\[0.2cm]
        Haushaltsgröße\\
        Anteil der Kinder\\[0.3cm]
        \scriptsize
        Welchen Einfluss hat die Haushaltsgröße,\\
        insbesondere der Anteil der Kinder,\\
        auf die Zugehörigkeit zur Armutsklasse?
    };

    \node[sub, below=of poverty] (struct)
    {
        \textbf{Strukturelles Risiko}\\[0.2cm]
        Berufsstatus\\
        Sektion / Tätigkeit\\[0.3cm]
        \scriptsize
        Inwieweit ist der Berufsstatus\\
        ein isolierter Prädiktor für Armut,\\
        wenn demografische Faktoren kontrolliert werden?
    };

    \node[sub, below right=of poverty] (social)
    {
        \textbf{Soziales Risiko}\\[0.2cm]
        Familienstand\\
        Witwer / Alleinstehende\\[0.3cm]
        \scriptsize
        Besteht ein signifikanter Zusammenhang\\
        zwischen Familienstand und extremer Armut?
    };

    % arrowssssssss
    \draw[arrow] (demo) -- (poverty);
    \draw[arrow] (struct) -- (poverty);
    \draw[arrow] (social) -- (poverty);

    \end{tikzpicture}

    \caption{Grafische Verfeinerung der Fragestellung gemäß Phase 1}
    \label{fig:fragestellung}
\end{figure}
